\section{Some complements}
\label{I.11}

We have already said that a connected \'etale covering
of an integral scheme is not necessarily integral.
Here are two examples of this fact.

a) Let $C$ be an algebraic curve with an ordinary double
point~$x$, $C^\prime$ its normalization, and $a$ and~$b$
the two points of~$C^\prime$ above~$x$. Let $C_i^\prime$
($i=1$, $2$) be two copies of~$C^\prime$, and $a_i$
and~$b_i$ the point of $C_i^\prime$ corresponding to $a$
\resp $b$. In the sum curve $C_1^\prime\amalg C_2^\prime$,
identify $a_1$ with~$b_2$, and $a_2$ with~$b_1$ (we leave
it to the reader to make the identification process
precise; it will be explained in Chap\ptbl VI of the
multiplodoque but, in the case of curves over an
algebraically closed field, is treated in Serre's book
on algebraic curves). One obtains a \emph{connected}
and \emph{reducible} curve $C^{\prime\prime}$ which is
an \'etale covering of degree~$2$ of~$C$. The reader
will verify that, in general, the connected ``Galois''
\'etale coverings $C^{\prime\prime}$ of~$C$ whose inverse
image $C^{\prime\prime}\times_C C^\prime$ is a
\emph{trivial} covering of~$C^\prime$ (\ie isomorphic
to the sum of a certain number of copies of~$C^\prime$)
are ``cyclic'' of degree~$n$, and that for every
integer~$n>0$ one can construct a connected cyclic
\'etale covering of degree~$n$. In the language of the
fundamental group, which will be developed later,
this means that the quotient of~$\pi_1(C)$ by the
closed normal subgroup generated by the image
of~$\pi_1(C^\prime)\to\pi_1(C)$ (the homomorphism
induced by the projection) is isomorphic to the
profinite completion of~$\ZZ$. More precisely, one
should be able to show that the fundamental group
of~$C$ is isomorphic to the (topological) free product
of the fundamental group of~$C^\prime$ with the
profinite completion of~$\ZZ$. Note that questions
of this kind gave rise to ``descent theory'' for schemes.

b) Let $A$ be a complete local integral domain. One
knows that its normalization $A^\prime$ is finite
% original p. 27
over~$A$ (Nagata), hence is a complete semilocal ring,
and therefore local since it is an integral domain.
Suppose that the residue field extension $L/k$ which
it defines is not purely inseparable (otherwise,
one will say that $A$ is geometrically unibranch;
\cf below). This is the case, for example, for the
ring $\RR[[s,t]]/(s^2+t^2)\RR[[s,t]]$, where $\RR$
is the field of real numbers. Let $k^\prime$ be a
finite Galois extension of~$k$ such that
$L\otimes_k k^\prime$ decomposes, and let~$B$ be
a finite \'etale algebra over~$A$ corresponding
to the residue field extension~$k^\prime$ (recall
that $B$ is essentially unique). Then
$B^\prime=A^\prime\otimes_A B$ over~$B$ has residue
algebra $L\otimes_k k^\prime$, which is not local;
hence $B^\prime$ is not a local ring, and therefore
(being complete) has zero divisors. Since it is
contained in the total ring of fractions of~$B$
(being free over~$A^\prime$, hence torsion-free
over~$A^\prime$, hence torsion-free over~$A$, and
therefore contained in
$B^\prime\otimes_A K=B^\prime_{(K)}
=A^\prime_{(K)}\otimes_K B_{(K)}=B_{(K)}$
since $A^\prime_{(K)}=K$), it follows that $B$ is
not an integral domain. In the case of the ring
$\RR[s,t]/(s^2+t^2)\RR[s,t]$, taking
$k^\prime/k=\CC/\RR$, one obtains for~$B$ the local
ring of two intersecting lines in the plane at
their point of intersection.

Note also that if there exists a connected \'etale
covering $X$ of an integral $Y$ which is not
irreducible, then every irreducible component of~$X$
gives an example of an unramified covering $X^\prime$
of~$Y$, dominating~$Y$, which is not \'etale over~$Y$.
In example~a), one thus finds that $C^\prime$ is
unramified over~$C$, without being \'etale at the
two points $a$ and~$b$ (as one also sees directly
by inspecting the completions of the local rings
of~$x$ and~$a$: from the ``formal'' point of view,
$C^\prime$ at the point~$a$ is identified with
a closed subscheme of~$C$ at the point~$x$, namely
one of the two ``branches'' of~$C$ through~$x$).

In a) and~b), one sees that the failure of the
conclusions of~\ref{I.9.5}\ptbl (i) and~(ii) is
directly linked to the fact that a point of~$Y$
``splits'' into \emph{distinct} points of the
normalization (in~b), the fact that the residue
field extension is not purely inseparable must
be interpreted geometrically in this way).
Precisely, we shall say that a local integral
domain~$A$ is \emph{geometrically unibranch} if
its normalization has only one maximal ideal,
with the corresponding residue field extension
purely inseparable; a point~$y$ of an integral
prescheme is said to be geometrically unibranch
if its local ring is so. Examples: a normal point,
an ordinary cusp of a curve, etc. It seems that
if $Y$ has a point which
% original p. 28
is not unibranch, there always exists a connected
\'etale covering of~$Y$ which is not irreducible;
at least this is what we have shown in case~b),
when $Y$ is the spectrum of a complete local ring.
One can show, on the other hand, that \emph{if
all points of~$Y$ are geometrically unibranch,
then every connected unramified $Y$-prescheme
dominating~$Y$ is \'etale} and irreducible.
The proof follows that of~\ref{I.9.5}, using the
following generalization of Theorem~\ref{I.8.3},
which will be proved later by means of descent
techniques\footnote{\Cf IX~\SourceRef{IX.4.10}{4.10}.
For a more direct proof, \cf EGA~IV 18.10.3,
using a variant of~\ref{I.9.5} for geometrically
unibranch local rings.}:

\emph{Let $Y^\prime\to Y$ be a finite, radicial,
surjective morphism (\ie what one might call
a ``universal homeomorphism''). Consider the
functor $X\mto X\times_Y Y^\prime=X^\prime$ from
$Y$-preschemes to $Y^\prime$-preschemes. This
functor induces an equivalence of the category
of \'etale $Y$-schemes with the category of
\'etale $Y^\prime$-schemes.} One can apply this
result, for example, when $Y^\prime$ is the
normalization of~$Y$, with $Y$ assumed unibranch
(and $Y^\prime$ finite over~$Y$, which holds in
all cases encountered in practice), or to a
$Y^{\prime\prime}$ ``sandwiched'' between~$Y$
and its normalization (which no longer needs
to be finite over~$Y$).
